% Copyright (c) 2026 Twin Vector Labs LLC.
% All rights reserved.
\documentclass[11pt]{article}

\usepackage[margin=1in]{geometry}
\usepackage{amsmath,amssymb,amsthm}
\usepackage{mathtools}
\usepackage{enumitem}
\usepackage{booktabs}
\usepackage{array}
\usepackage[colorlinks=true,linkcolor=blue,urlcolor=blue]{hyperref}

\newtheorem{principle}{Principle}
\theoremstyle{definition}
\newtheorem{definition}{Definition}
\newtheorem{remark}{Remark}

\newcommand{\C}{\mathbb{C}}
\newcommand{\R}{\mathbb{R}}
\newcommand{\Z}{\mathbb{Z}}
\newcommand{\rep}[1]{\mathbf{#1}}
\newcommand{\triv}{\mathbf{1}}
\newcommand{\abs}[1]{\left\lvert #1 \right\rvert}
\newcommand{\norm}[1]{\left\lVert #1 \right\rVert}
\newcommand{\sigmacol}{\sigma}
\newcommand{\Pin}{P_{\mathrm{in}}}
\newcommand{\Pout}{P_{\mathrm{out}}}
\newcommand{\dd}{\mathrm{d}}
\DeclareMathOperator{\tr}{tr}
\DeclareMathOperator{\sgn}{sgn}

\title{\bfseries The Flavor and Charge Sectors:\\
a spatio-temporal and Dirac--K\"ahler program for the emergent proton's
electroweak content}
\author{Tessera --- cobordism subsystem (epic \#410: \#411--\#418)}
\date{}

\begin{document}
\maketitle

\begin{abstract}
The cobordism subsystem now carries the \emph{color-confinement skeleton} of the
proton: on the symmetric trivalent junction $W_{ABC}$ a color-symmetric quark input
transports to the net-color-zero singlet $\langle(1,\omega,\omega^2)\rangle$ with
overlap $0.999$, confinement ($\sigmacol\to0$) is a Stokes conservation law at the
junction, and the location of the singlet is forced by the geometry's $\Z_3$ symmetry
rather than imposed. What the construction does \emph{not} carry is the electroweak
content: electric charge, the proton/neutron (flavor) distinction, and the
spinor/current structure. The earlier conclusion (\#405) was that these need a
\emph{richer construction} --- a flavor register parallel to color, plus spinor
fields. This document lays out a different program, the one we are now building and
testing: find the missing content by \textbf{splitting the behavior of the existing
register} spatio-temporally and across form-degree, and by removing two arbitrary
choices that currently break the geometry's symmetry. Concretely: (i) source simplex
orientation/sign from the standard simplicial formula, never a vertex-label sort; (ii)
replace the asymmetric prism extrusion with a symmetric simplicial interior; (iii) read
electric charge as a temporal-sector Gauss-law holonomy; (iv) realize the spinor/current
content as the Dirac--K\"ahler factor of the Hodge operator; and (v) seek flavor in the
spatio-temporal split or the Dirac--K\"ahler multiplicity, \emph{not} a parallel
register. We define each sector, its observable, and its falsifiable test; we fix a set
of faithfulness invariants ($G_1$--$G_8$) that no subtask may regress; and we give the
dependency structure that says what can be built in parallel. This is a \emph{plan},
honest about what is established versus hypothesis; the measurements live in the
subtasks.
\end{abstract}

\tableofcontents

\section{Where we are, and what is missing}
\label{sec:where}

\paragraph{Established.} Matter is not imposed in this subsystem. Color states are
pinned on the boundary $\partial W$ of a Lorentzian Regge cobordism $W$, the interior
edge lengths relax to a stationary point of the dual Regge action, and the emergent
result is \emph{read off} the relaxed geometry. Color lives in the homology of a holed
register surface; the proton is the trivalent junction $W_{ABC}$ --- one connected
$2$-frequency geodesic icosahedron ($S^2$, $42$ vertices) minus twelve vertex-disjoint
hole triangles in four windows of three ($A,B,C$ inputs and $R$ the emergent result),
extruded $\times I$. Three facts are in hand:
\begin{itemize}[nosep]
  \item \textbf{The emergent color singlet.} With the windows placed as one
  tetrahedral $A_4$ orbit, the input$\to$result transport $M$ intertwines the color
  $\Z_3$, $M\Pin = \Pout M$, and the color-symmetric ($\omega$-representation) input
  transports to the singlet line with overlap $0.999$.
  \item \textbf{Confinement as conservation.} The global Stokes relation
  $\sigmacol_R = -(\sigmacol_A+\sigmacol_B+\sigmacol_C)$ holds exactly
  ($\sim 10^{-14}$ on the uniform metric); a net color charge has no realizable
  harmonic representative.
  \item \textbf{The right dimensionless ballpark.} Read emergently,
  $\sigmacol\to0$, and $r\!\cdot\!m\approx 6.9$ against the physical proton's $4.0$.
\end{itemize}

\paragraph{The boundary (\#405).} Every electroweak quantum number needs structure the
\emph{color-only} geometry does not carry. The three windows are identical color
registers, so nothing distinguishes the proton (uud, $+1$) from the neutron (udd, $0$):
that is \emph{flavor}. Electric charge is flavor-dependent (u${:}+\tfrac23$,
d${:}-\tfrac13$). Spin needs spinor content; the construction is a geometry of scalar/
complex cochains. The natural ``net orientation as charge'' read reduces to the (zero)
color flux or a particle/antiparticle sign --- not the flavor-dependent electric charge.

\paragraph{The fork.} The earlier write-ups proposed the minimal extension as a
\emph{richer construction}: a flavor register parallel to color, and spinor fields. We
take a different road, for two reasons. First, a second hole register is geometrically
expensive and competes with color for the surface's vertex capacity, and if it shares
the color holes it is no longer independent --- it double-counts. Second, and more
importantly, the missing content has a natural home in structure the construction
\emph{already has} but does not yet use: the Lorentzian (temporal/spatial) split of its
cochains, and the higher-degree pieces of the cochain complex. The program of this
document is to use that structure.

\section{The reframe}
\label{sec:reframe}

\begin{principle}[Split the register, do not multiply it]
\label{prin:split}
The proton's electroweak content is sought by splitting the behavior of the
\emph{existing} register --- across causal type (temporal vs.\ spatial) and across
form-degree (the inhomogeneous cochain) --- rather than by adding parallel hole
registers. Independence between sectors is to be supplied by orthogonality in the
spectral/causal decomposition, not by disjoint topology.
\end{principle}

Three observations make this concrete.

\paragraph{Charge is a temporal-sector object.} The electromagnetic field strength is a
$2$-cochain $F\in\Omega^2$. On a Lorentzian complex it splits by the causal type of the
plaquette: the \emph{electric} field is the part with a timelike leg, the \emph{magnetic}
field is the purely spacelike part. Electric charge is then Gauss's law on the temporal
sector, $\dd{\star F}=\star j$ with $\oint \star F = Q$ --- a genuine gauged-$U(1)$
holonomy, quantized and metric-robust, not a hand-chosen weight. The quark charges are
multiples of $\tfrac13$, and the $\tfrac13$-quanta are tied to \emph{triality} (the
color $\Z_3$ center the construction already carries) through Gell-Mann--Nishijima
$Q=I_3+\tfrac12 Y$.

\paragraph{Spin/current is the Dirac--K\"ahler factor of the operator.} The inhomogeneous
cochain $\Omega^\bullet=\bigoplus_k\Omega^k$ is the Dirac--K\"ahler field; the operator
$(\dd+\delta)$, whose square is the Hodge Laplacian $L=(\dd+\delta)^2$, is the discrete
Dirac operator, with gamma matrices realized as the Clifford action of edges
($1$-cochains). In four dimensions $\dim\Omega^\bullet = 16 = 4\times4$ and the operator
algebra is the Dirac algebra $\C l(1,3)\otimes\C$: a genuine $4\times4$ gamma structure,
\emph{factored out of the Hodge operator the construction already builds}. Its conserved
current's time component $j^0$ is the charge density. This is also the principled
distinguished frame the deferred operator read-out lacks (\S\ref{sec:dk}).

\paragraph{The geometry must be honest.} Two arbitrary choices currently break the
symmetry the physics needs. The simplex orientation/triangulation is in places decided
by a vertex-label \texttt{std::sort}, not the standard simplicial orientation; and the
$\times I$ prism extrusion is asymmetric --- it is the documented source of the
$\norm{M\Pin-\Pout M}/\norm{M}\approx 4.3\times10^{-2}$ intertwining residual that pins
the singlet at $0.999$ instead of $1$. Both must be removed before the spectral/causal
split can be trusted.

\medskip
\noindent The remainder of this document develops these into a buildable program. Each
section names the subtask issue that carries it and the falsifiable test that keeps it
faithful.

\section{Geometric foundations}
\label{sec:foundations}

The two foundations remove arbitrariness; everything downstream inherits the symmetry
they restore.

\subsection{Standard simplex orientation, not label sort (\#412)}
\label{sec:orientation}

The standard orientation of a $k$-simplex is an equivalence class of vertex orderings
under \emph{even} permutations; the sign of a given ordering is the parity of the
permutation relative to a \emph{consistent induced (global) orientation}. Choosing the
sign by sorting \emph{vertex labels} makes the result depend on the arbitrary numbering
of the complex, which is not a symmetry of the geometry. The icosahedral automorphism
$g$ that realizes the color $\Z_3$ permutes vertices but not labels; a label-sort sign
is therefore not $g$-equivariant, and the symmetry leaks into the read-out.

\begin{principle}[Orientation is geometric]
\label{prin:orient}
Every simplex orientation/sign --- in the boundary operator, the period read-out
covector, and any triangulation diagonal --- is sourced from the standard simplicial
orientation formula tied to the manifold's induced orientation
(\texttt{SimplexOrientation::}\allowbreak\texttt{orientationOf}). Sorting survives only as a sign-neutral
canonical storage key, never as a sign.
\end{principle}

\paragraph{Faithfulness.} The test of correctness is \emph{relabeling-invariance}:
applying a random permutation to the base-surface vertex ids and rebuilding $W_{ABC}$
must leave the color charge $\sigmacol$, the singlet overlap, and the intertwining
residual unchanged up to one consistent global sign shared by all windows. This is the
invariant $G_6$ of \S\ref{sec:faithful}; \#412 owns it.

\subsection{A symmetric simplicial interior (\#413)}
\label{sec:interior}

The cobordism interior is currently the Freudenthal ``staircase'': each base cell is cut
into simplices along a diagonal selected by sorting vertex ids. That cut is
label-dependent and not $g$-invariant, and it --- not the metric --- is the cause of the
$4.3\times10^{-2}$ residual: replacing the uniform seed by a $g$-invariant
non-degenerate metric leaves the residual essentially unchanged. The fix is a
\emph{symmetric apex stacking} --- a continuation of the primal $1$-skeleton, \emph{not a
prism}: each base triangle at $t=0$ cones up to a single \emph{face-apex} vertex $f$ at
$t=1$ (a $(3,1)$ tetrahedron), and that apex is shared with the inverted triangle at $t=2$
(a $(1,3)$ tetrahedron). The gap between two adjacent apex-columns over a shared base edge
is an octahedron split along the \emph{canonical dual edge} $f_1 f_2$ (the two
face-centres) into four tetrahedra --- so \emph{no vertex-label sort chooses any diagonal
anywhere}. It is label-independent and yields genuine vertical worldline (temporal) edges.

\begin{remark}[Connectivity is CDT, the metric stays Regge]
\label{rem:cdt}
The alternating stacking is the Ambj\o rn--Loll Causal Dynamical Triangulations
\emph{connectivity}; the codebase already carries that orientation convention. But we do
\emph{not} import CDT's preferred time foliation: the edge lengths stay free, the causal
type emergent, the action stationary. Connectivity and the metric are independent
choices; we take the symmetric connectivity and keep the Regge metric.
\end{remark}

This foundation does double duty. It removes the residual (driving the singlet overlap to
$1$), and its worldline edges are the clean temporal edges the charge sector needs. The
octahedral cells are split along the canonical dual edge $f_1 f_2$, determined by the two
incident faces rather than by vertex ids, so the split is itself $g$-equivariant --- there
is no residual diagonal choice left to document.

\paragraph{Faithfulness (measured).} A prototype of the apex stacking on the real
$42$-vertex base drives the residual $\norm{M\Pin-\Pout M}/\norm{M}$ from the prism's
$4.3\times10^{-2}$ to $4.5\times10^{-14}$ (machine zero) and the singlet overlap from
$0.999$ to $1.000000$, with $b_1=11$, a valid manifold (\texttt{dualComplexValid}), and
charge conservation exact --- the transport is \emph{exactly} equivariant. \#413 owns
``singlet $\to1$''.

\section{The charge sector}
\label{sec:charge}

\subsection{E and B from the temporal/spatial split (\#417)}
\label{sec:eb}

Let $A\in\Omega^1$ be a $U(1)$ connection and $F=\dd A\in\Omega^2$ its curvature. On the
Lorentzian complex each plaquette (a $2$-cell, dually a hinge) is classified by the
causal type of its edges:
\begin{equation}
  F = \underbrace{F^{(E)}}_{\text{plaquettes with a timelike leg}}
      \;\oplus\;
      \underbrace{F^{(B)}}_{\text{purely spacelike plaquettes}} ,
\end{equation}
the discrete analogue of $E_i=F_{0i}$, $B_{ij}=F_{ij}$. The split is read off the
existing geometry by causal type --- no new register, no new holes. The reader is a pure
observable: it must leave every color observable of \S\ref{sec:where} bit-for-bit
unchanged.

\paragraph{Faithfulness.} On a purely spacelike (static) seed $\norm{E}=0$ and $B$ carries
the structure; introducing timelike flux makes $\norm{E}>0$; and $F=\dd A$ is
gauge-invariant ($A\to A+\dd\chi$ leaves $E,B$ unchanged). A wrong causal partition leaks
$B$ into $E$ and the static-seed test catches it.

\subsection{Charge as a Gauss-law holonomy (\#411)}
\label{sec:gauss}

Electric charge is the discrete Gauss law: the flux of $E$ through a closed surface
enclosing the quark windows,
\begin{equation}
  Q \;=\; \oint_{\partial V} \star F \;=\; \int_V \star j ,
  \label{eq:gauss}
\end{equation}
the time component of the conserved current. The point is what kind of object this is.
The previously-considered ``flavor-weighted covector'' $(\tfrac23,\tfrac23,-\tfrac13)$ is
not $A_4$-equivariant and is not the topologically protected functional: it would degrade
under a non-symmetric seed the way charge conservation degrades from $10^{-14}$ to
$\sim10^{-2}$ under jitter. The Gauss-law holonomy \eqref{eq:gauss}, by contrast, is a
genuine gauged-$U(1)$ invariant --- quantized and metric-robust.

\paragraph{Triality and fractional charge.} Quark electric charges are multiples of
$\tfrac13$, and in the Standard Model the fractional part is tied to color through the
center $\Z_3$ (triality) via $Q=I_3+\tfrac12 Y$, with $\tfrac12 Y$ carrying baryon
number/triality. The construction already carries triality (the color $\Z_3$ is
$\Pout$). This sector investigates whether the $\tfrac13$-quantum of the Gauss-law charge
\emph{is} that triality --- i.e.\ whether fractional electric charge is forced by the
same $\Z_3$ that forces confinement.

\paragraph{Faithfulness.} Over an ensemble of metric jitters, $\mathrm{std}(Q)$ stays
tiny while the covector baseline drifts by orders of magnitude (the ratio is the figure
of merit); $Q$ lands on the third-integer lattice; and on the color singlet $Q$ equals
the expected neutral total.

\section{The spin/current sector: Dirac--K\"ahler}
\label{sec:dk}

\subsection{The operator and its gamma structure (\#415)}

The inhomogeneous cochain $\Omega^\bullet=\bigoplus_{k}\Omega^k$ is the discrete
Dirac--K\"ahler field. With the coboundary $\dd$ and its adjoint (codifferential)
$\delta$, the Dirac--K\"ahler operator is
\begin{equation}
  D = \dd+\delta, \qquad D^2=(\dd+\delta)^2=\dd\delta+\delta\dd = L,
  \label{eq:dk}
\end{equation}
the Hodge Laplacian. Thus $D$ is a \emph{square root} of the operator the construction
already assembles. Clifford multiplication by a $1$-cochain $e$ acts on $\Omega^\bullet$
by $e\vee\alpha = e\wedge\alpha + \iota_e\alpha$ (an edge raises and lowers form degree),
and these satisfy the Clifford relations
\begin{equation}
  \{\gamma_a,\gamma_b\}=2\,\eta_{ab}.
\end{equation}
In four dimensions $\Omega^\bullet\cong\C l(1,3)\otimes\C\cong M_4(\C)$, the $4\times4$
Dirac algebra: the gamma matrices are \emph{factored out of the operator}, exactly the
intuition that the spin matrix is a factor of the (Choi/operator) object of the right
dimension.

\paragraph{The distinguished frame.} The operator read-out
$U=\mathrm{unvec}(\ker L_1(W-\partial W))$ is presently deferred because the kernel basis
is fixed only up to an $O(d^2-1)$ rotation --- there is no canonical frame. The
Dirac--K\"ahler structure supplies one: the Clifford generators are a distinguished basis
with distinct eigenstructure, exactly as $\Pout$'s distinct eigenvalues
$\{1,\omega,\omega^2\}$ gave color its canonical frame. Pinning the operator frame to the
gamma generators is the same de-framing trick, transplanted from the state sector to the
operator sector.

\subsection{The current as charge density, and the multiplicity}

The Dirac--K\"ahler current $j$ has time component $j^0$ equal to the charge density;
this must agree with the Gauss-law charge of \S\ref{sec:gauss}, tying the operator sector
to the temporal sector. The construction has a known feature: the Dirac--K\"ahler field
in $d$ dimensions carries $2^{d/2}$ Dirac spinors --- a four-fold multiplicity in 4D (the
``doubling''). On the lattice this is usually a nuisance; here it is a candidate
\emph{flavor/taste} index, a built-in internal multiplicity that may furnish flavor
without a second register (\S\ref{sec:flavor}). We record it honestly as a hypothesis,
not a result.

\begin{remark}[Spin is an operator property; color/flavor are state properties]
This sector placement is not arbitrary. Spin is how a state transforms under rotation ---
an \emph{operator}. Color (and, below, flavor) are properties of the boundary state's
periods. The qudit form of the singlet obstruction already hints at this: the
qubit ($d{=}2$) alternating invariant $\varepsilon_{ij}$ \emph{is} the spin singlet,
bipartite, living in the $d{=}2$ operator sector, while the qutrit color singlet
$\varepsilon_{ijk}$ is irreducibly tripartite in the state sector.
\end{remark}

\section{The flavor sector}
\label{sec:flavor}

Proton (uud) and neutron (udd) differ only by flavor; the color-only geometry cannot tell
them apart. By Principle~\ref{prin:split} we seek flavor as a \emph{split of the existing
register}, not a parallel hole register. Two candidate mechanisms are on the table:
\begin{enumerate}[nosep,label=(\roman*)]
  \item the spatio-temporal split --- does the temporal versus spatial behavior of the
  register furnish a binary, u/d-like label?
  \item the Dirac--K\"ahler multiplicity (\S\ref{sec:dk}) --- does its four-fold
  taste index carry flavor?
\end{enumerate}
The success criterion is concrete and shared with the charge sector: a geometric handle
that yields total electric charge $+1$ for uud and $0$ for udd.

\begin{remark}[Why ``flavor as a second nonzero spectral band'' is rejected]
A tempting reading of ``flavor as a second harmonic'' is the first \emph{nonzero}
eigenvalue band of $L_1$. This is unsound. A period/holonomy is defined only on
\emph{closed} forms (the kernel); a $\lambda\neq0$ eigenform has $\dd\psi\neq0$, so its
``period'' is representative-dependent and not Stokes-conserved --- and the charge sum of
\S\ref{sec:charge} type-checks only on the kernel. A generic nonzero band is also
multiplicity one, so it cannot host an isospin doublet. Flavor must be either a second
\emph{kernel} structure or, as proposed here, the spatio-temporal / Dirac--K\"ahler split
--- not a free-floating band.
\end{remark}

\paragraph{Faithfulness.} The discriminator $Q(\text{uud})=+1$, $Q(\text{udd})=0$ is the
forward test. The decisive \emph{structural} guard is that $b_1$ stays $11$: if a new
independent cycle appears, a register has been smuggled in, contradicting
Principle~\ref{prin:split}. Proton and neutron must differ only in the flavor label, not
in color content ($\sigmacol\to0$ for both).

\section{The binding structure: a bipartite revisit (\#416)}
\label{sec:bipartite}

The proton was placed on a trivalent junction because the color singlet
$\varepsilon_{ijk}$ is absent from $\rep3\otimes\rep3=\rep6\oplus\bar{\rep3}$ and first
appears in $\rep3^{\otimes3}$. That theorem rules out a \emph{generic} bipartite merge:
the relaxed two-quark result carries $\rep6\oplus\bar{\rep3}$ and the symmetric $\rep6$
contaminates. It does \emph{not} rule out two bipartite steps with a projection: if the
first merge projects onto the antisymmetric diquark $\bar{\rep3}=\varepsilon q q$, the
second ($\bar{\rep3}\otimes\rep3=\triv\oplus\rep8$) contains the singlet.

Two questions follow, both testable with machinery that already exists --- the
\emph{bipartite} merge's prism extrusion accepts a cumulative \emph{twist} permutation,
i.e.\ an orientation-reversing tube that is a geometric antisymmetrizer. (The
\emph{tripartite} $W_{ABC}$ interior is now the symmetric apex stacking of
\S\ref{sec:interior}, with no label-sort diagonal; the bipartite path here still uses the
prism, and a symmetric analog of the twist would be its eventual replacement.)
\begin{itemize}[nosep]
  \item Does a twisted bipartite tube project the emergent diquark onto $\bar{\rep3}$
  ($\abs{\sigmacol}\approx1$, the antisymmetric channel) rather than the generic average
  ($\abs{\sigmacol}\approx0.67$)?
  \item The deeper obstruction is that the colored diquark intermediate is non-carriable
  (confinement): only $\sigmacol=0$ combinations carry as free boundary states. A
  \emph{connected} two-vertex ``path'' geometry, where the diquark is internal and never
  a free boundary, evades this. Does such a path reach the singlet, and is it genuinely
  distinct from the symmetric star junction?
\end{itemize}
This sector is included because it bears directly on whether the proton's binding (and
thus its baryon number) is irreducibly a star or admits a path realization --- a
structural question underneath flavor and charge.

\paragraph{Faithfulness.} The twisted diquark's $\abs{\sigmacol}$ must clearly separate
from the generic value; the path geometry's singlet overlap is pinned to a recorded
verdict; the intermediate's carriability residual is reported so a ``separated two-step''
claim cannot quietly rely on a non-carriable free diquark.

\section{The dimensional question (\#418, \#429)}
\label{sec:dim}

The full structure of this program wants the full dimension. The electromagnetic field
strength has its full $\vec E,\vec B$ content, and the Dirac algebra its full $4\times4$
gammas, in $3+1$ dimensions --- a genuine $S^3$ spatial slice extruded in time. The
discipline here is the program's discipline everywhere: we do \emph{not} reduce dimension
as a shortcut. Tessera is fully emergent, and dropping to $2+1$ D (an $S^2$ slice)
\emph{changes the physics} --- the field strength loses a component ($B$ becomes a scalar)
and the Clifford algebra collapses to $\C l(2,1)$ ($2\times2$ gammas). A result measured
in a reduced sector is a result about a different theory, not a faithful first cut of this
one. There is therefore no ``explore it in $2+1$ D first'': the Gauss-law charge, the
binary flavor label, the full $\vec E,\vec B$, and the $4\times4$ Dirac--K\"ahler are all
to be measured in the full dimension.

The apparent obstacle --- that the symmetric apex interior of \S\ref{sec:interior} is
currently triangle-only --- is an implementation gap, not a dimensional wall. The interior
connectivity is \emph{dimension-generic}, defined entirely from the base layer's
coface/dual structure: cone each top $d$-simplex to an apex vertex, and let the
apex-connectivity be \emph{mirrored} across shared faces. A facet (a $(d-1)$-simplex)
shared by two cofaces has those cofaces' apex-edges reflected across it; a lower face
shared by $k$ cofaces mirrors analogously. In $2$D this mirroring \emph{is} the dual-edge
octahedron split of \S\ref{sec:interior} (a shared edge's two triangle-apexes reflected
through it); in $3$D it cones tetrahedra to cell-apexes across shared triangles; and so on
in every dimension --- no vertex-label sort chooses any diagonal at any dimension. The
current \texttt{Spacetime::symmetricStackCells} triangle-only special case (the
\texttt{t.size()\,!=\,3 -> continue} guard) is the temporary gap, not a barrier. The full
dimension is reached by \emph{generalizing} this construction, which is the content of
\#429, not by reducing to a sector.

\#429 (\emph{iterated apex-reflection cobordism, dimension-generic, and emergent twists})
carries the generalization: lift \texttt{symmetricStackCells} to the coface-mirroring rule
above, add an \texttt{nApexSlices} parameter (default $1$, today's single stack), build an
$S^3$ over a tetrahedral base, and measure the emergent twist of \S\ref{sec:twist}. \#418
remains the go/no-go scoping --- what the Hodge/Regge/relaxation machinery costs over a
triangulated $S^3$, and a minimal $S^3$ smoke test --- but the recommendation it returns
is which observables to \emph{build} in the full dimension, never which to retreat from it.

\subsection{Iterated apex-reflection and the emergent twist (\#429)}
\label{sec:twist}

The symmetric apex interior has a reading that the prism never had. Each apex is a centre
of \emph{point-reflection}: the next slice is made by inverting an apex's neighbour edges
through it (and adding a ``cap'' to close the cell). A single slice (\texttt{nApexSlices}
$=1$) is today's construction; iterating reflect-and-cap over many slices stacks the
cobordism in time. Because a point-reflection is an orientation-reversing involution and
the \emph{composition of two reflections is a rotation}, an iterated stack of apex
reflections can accumulate a net rotation --- a \emph{screw}/monodromy --- along the time
direction. That accumulated rotation is a \textbf{twist}.

The twist is \emph{emergent}: it is read off the relaxed geometry of the iterated stack,
not supplied by hand. This is the crucial contrast with the bipartite tube's
\texttt{prismCells} twist of \S\ref{sec:bipartite}, which is a hand-chosen permutation
inserted into the extrusion. A nontrivial emergent twist makes $W$ a \emph{mapping torus}
--- a twisted product, not a plain $\Sigma\times I$ --- with the twist as its gluing
monodromy. The hypothesis to test in \#429 is whether such a twist appears and what it
means:
\begin{itemize}[nosep]
  \item a $\Z_3$ twist would be the color singlet's $\omega$-phases (triality realized
  geometrically along time);
  \item a $2\pi\!\mapsto\!-1$ twist would be the spinor sign --- genuine spin, the central
  $\Z_2$ of $\mathrm{Spin}\to SO$ flagged as a risk in \S\ref{sec:risks};
  \item a $U(1)$ twist would be a gauge holonomy --- the charge of \S\ref{sec:charge} as a
  monodromy.
\end{itemize}
Each reading is a falsifiable target: the screw angle of the relaxed stack is measured and
matched (or not) to one of these. As with every observable in this program, the twist is
something we \emph{read}, not impose; \#429 prototypes it over the iterated apex stack and
its $S^3$ generalization.

\section{Faithfulness and anti-drift}
\label{sec:faithful}

This is a multi-step program over one shared geometry; later subtasks must not silently
regress earlier results. We fix the \emph{golden} invariants --- the proton facts already
established --- as a shared regression guard, and require every subtask's tests to keep
them.

\begin{center}
\small
\setlength{\tabcolsep}{5pt}
\begin{tabular}{@{}clp{0.62\linewidth}@{}}
\toprule
 & invariant & assertion (the ``do not drift'' guard)\\
\midrule
$G_1$ & proton present & singlet overlap $0.999$ (prism) $\to 1.000000$ (symmetric apex interior, \#413)\\
$G_2$ & confinement & color charge $\sigmacol\to0$ ($\abs{\sigmacol}\lesssim0.06$) for the symmetric input\\
$G_3$ & charge conservation & $\sigmacol_R=-(\sigmacol_A+\sigmacol_B+\sigmacol_C)$ exact ($\sim10^{-14}$, uniform metric)\\
$G_4$ & topology & $b_1=11$ and the dual complex is a valid manifold (no weld)\\
$G_5$ & color $\Z_3$ intact & $\Pout$ eigenvalues $\{1,\omega,\omega^2\}$\\
$G_6$ & relabeling-invariance & permuting vertex ids leaves $\sigmacol$/overlap/residual invariant (owned by \#412)\\
$G_7$ & determinism & fixed seed $\Rightarrow$ reproducible result; no wall-clock/random seeding\\
$G_8$ & emergent-first & observables read off the relaxed geometry, never hand-placed\\
\bottomrule
\end{tabular}
\end{center}

\noindent Each subtask carries a \emph{Faithfulness \& anti-drift tests} section that
pins its own forward test (expected values, fixed seed) and names the subset of
$G_1$--$G_8$ it must keep. They share one module,
\texttt{tests/cobordism/test\_epic410\_invariants.py}; every pull request in this track
must keep it green. The discipline is sector-specific: an observable-only sector
(\S\ref{sec:eb}) must leave $G_1$--$G_5$ bit-for-bit unchanged; a geometry sector
(\S\ref{sec:interior}) must preserve $G_3$/$G_4$; the flavor sector
(\S\ref{sec:flavor}) must keep $b_1$ unchanged so flavor is not smuggled in as holes.

\section{Dependency structure and sequencing}
\label{sec:deps}

The subtasks form a small dependency graph; ``$A\to B$'' means $B$ needs $A$.

\begin{center}
\ttfamily\small
\begin{tabular}{@{}l@{}}
\#412 orientation --> \#413 interior --+--> \#417 E/B --> \#411 charge --+\\
\hphantom{\#412 orientation --> \#413 interior } +--> \#415 Dirac-Kahler ---+--> \#414 flavor\\
\#416 bipartite twist\hphantom{xxxx}(soft: \#412 clean signs)\\
\#418 S\textasciicircum3 spike\hphantom{xxxxxxx}(investigation; informs \#417/\#411/\#415)\\
\end{tabular}
\end{center}

\noindent The parallelization waves:
\begin{itemize}[nosep]
  \item \textbf{Wave 0} (start in parallel): \#412 (orientation, the root), \#416
  (bipartite, independent), \#418 ($S^3$ spike, early so it can scope the rest).
  \item \textbf{Wave 1}: \#413 (symmetric interior), after \#412.
  \item \textbf{Wave 2} (parallel): \#417 (E/B) and \#415 (Dirac--K\"ahler), after \#413.
  \item \textbf{Wave 3}: \#411 (Gauss-law charge), after \#417.
  \item \textbf{Wave 4}: \#414 (flavor, the proton/neutron discriminator), after \#411
  and \#415.
\end{itemize}
The edges are forced: \#413 needs \#412 so the symmetric tetrahedralization does not
re-import a label-sort diagonal; \#417 needs \#413's worldline edges for a clean causal
split; \#411 needs \#417's $E$; \#414's success metric is the \#411 charge and its second
mechanism is the \#415 multiplicity.

\section{Open questions and honest risks}
\label{sec:risks}

This program is a plan; several of its load-bearing claims are hypotheses that the
subtasks must confirm or kill.

\begin{itemize}[nosep]
  \item \textbf{Spin-statistics / total antisymmetrization (the deepest gap).} A proton
  is the totally-antisymmetric $SU(6)$ qqq state: color (already antisymmetric, the
  $\varepsilon_{ijk}$ line) \emph{forces} flavor $\times$ spin $\times$ spatial to be
  symmetric and entangled (the $\mathbf{56}$-plet). $\lvert p\!\uparrow\rangle$ does not
  factorize into flavor $\otimes$ spin. Three sectors treated independently --- color
  (state), flavor (state), spin (operator) --- supply no mechanism for that cross-sector
  entanglement; placing spin in the operator sector and flavor in the state sector
  arguably forbids it. This is the program's largest unresolved question and must be held
  in view throughout.
  \item \textbf{Isospin is $SU(2)$, not a $\Z_2$ label.} A faithful flavor sector needs
  the isospin $SU(2)$ to \emph{act} (the ladder $p\leftrightarrow n$), as the color
  $\Z_3$ acts via the $\Pout$ intertwiner. A static binary u/d label is weaker; whether
  the spatio-temporal or Dirac--K\"ahler split furnishes a genuine doublet with a ladder
  is open.
  \item \textbf{The Dirac--K\"ahler doubling.} The four-fold multiplicity is a feature
  only if it can be \emph{interpreted}; it may instead be a discretization artifact. Its
  status as flavor/taste is a hypothesis.
  \item \textbf{Spinor double cover.} Genuine spin-$\tfrac12$ requires the central $\Z_2$
  of $\mathrm{Spin}\to SO$, distinct from the determinant/orientation $\Z_2$; the
  Dirac--K\"ahler construction must be checked to realize the former (a $2\pi$ rotation
  $\mapsto -1$), not merely the latter.
  \item \textbf{Metric dependence.} The temporal/spatial split and the nonzero spectral
  content are metric-dependent and move under relaxation, unlike the topological color
  kernel. This is physically apt (electric/flavor structure is not exactly conserved) but
  demands care that observables are read at a well-defined point and are not drifting
  artifacts.
\end{itemize}

\section{Scope and relationship to the existing program}
\label{sec:scope}

This is its own exploration track (epic \#410), independent of the proton epic's (\#380)
sequencing and deferrals. It does not modify the color result; it builds outward from it.
The geometric foundations (\S\ref{sec:foundations}) also improve the color result
directly (the symmetric interior drives the singlet overlap to $1$), so they are worth
doing regardless of how far the electroweak sectors reach. The honest framing is the same
as the color program's: we read what the geometry produces, we do not impose the answer,
and we let the falsifiable tests of \S\ref{sec:faithful} arbitrate.

\end{document}
